\documentclass[12pt,a4paper]{br-lex}

\usepackage[utf8]{inputenc}
\usepackage[T1]{fontenc}
\usepackage[brazil]{babel}
\usepackage{microtype}
\usepackage{booktabs}
\usepackage[hidelinks]{hyperref}

\title{Diretrizes para a interoperabilidade de dados públicos}
\author{Your Name}
\date{\today}

\begin{document}
\maketitle

\begin{center}
\begin{minipage}{0.88\textwidth}
\small
Este texto demonstra uma estrutura compacta para uma proposta normativa em
português: ementa, preâmbulo, capítulos, artigos e disposições finais. O
conteúdo é fictício e serve apenas como ponto de partida editorial.
\end{minipage}
\end{center}

\section*{Ementa}
Estabelece princípios para a publicação de conjuntos de dados públicos em
formatos interoperáveis e define responsabilidades mínimas de documentação.

\section{Disposições preliminares}
\subsection{Objeto e âmbito de aplicação}
\textbf{Art. 1º} Esta proposta estabelece requisitos para a publicação,
documentação e atualização de dados públicos destinados ao uso por pessoas,
organizações e sistemas automatizados.

\textbf{Parágrafo único.} As disposições aplicam-se aos conjuntos de dados
produzidos ou custodiados pelo órgão responsável, respeitadas as hipóteses de
sigilo e proteção de dados pessoais.

\section{Princípios de publicação}
\textbf{Art. 2º} A publicação observará os seguintes princípios:
\begin{enumerate}
  \item descrição clara da origem, da finalidade e da periodicidade de
    atualização;
  \item uso de formatos abertos e documentação suficiente para o reuso; e
  \item indicação de limitações conhecidas, inclusive lacunas e alterações
    metodológicas.
\end{enumerate}

\textbf{Art. 3º} Cada versão publicada deverá conservar um identificador e a
data de sua disponibilização, de modo a permitir auditoria e reprodução das
análises realizadas a partir dela.

\section{Responsabilidades}
\textbf{Art. 4º} Compete à unidade responsável:
\begin{itemize}
  \item manter um catálogo atualizado;
  \item responder a pedidos de esclarecimento técnico; e
  \item corrigir erros materiais sem apagar o histórico de versões.
\end{itemize}

\section{Disposições finais}
\textbf{Art. 5º} Esta proposta entra em vigor na data de sua publicação.

\end{document}
